\documentclass[10pt, letterpaper]{article}

% Packages:
\usepackage[
    ignoreheadfoot, % set margins without considering header and footer
    top=1 cm, % seperation between body and page edge from the top
    bottom=1 cm, % seperation between body and page edge from the bottom
    left=1 cm, % seperation between body and page edge from the left
    right=1 cm, % seperation between body and page edge from the right
    footskip=1.0 cm, % seperation between body and footer
    % showframe % for debugging 
]{geometry} % for adjusting page geometry
\usepackage{titlesec} % for customizing section titles
\usepackage{tabularx} % for making tables with fixed width columns
\usepackage{array} % tabularx requires this
\usepackage[dvipsnames]{xcolor} % for coloring text
\definecolor{primaryColor}{RGB}{0, 0, 0} % define primary color
\usepackage{enumitem} % for customizing lists
\usepackage{fontawesome5} % for using icons
\usepackage{amsmath} % for math
\usepackage[
    pdftitle={Karteek's CV},
    pdfauthor={Karteek Gandiboyina},
    pdfcreator={LaTeX with RenderCV},
    colorlinks=true,
    urlcolor=primaryColor
]{hyperref} % for links, metadata and bookmarks
\usepackage[pscoord]{eso-pic} % for floating text on the page
\usepackage{calc} % for calculating lengths
\usepackage{bookmark} % for bookmarks
\usepackage{lastpage} % for getting the total number of pages
\usepackage{changepage} % for one column entries (adjustwidth environment)
\usepackage{paracol} % for two and three column entries
\usepackage{ifthen} % for conditional statements
\usepackage{needspace} % for avoiding page brake right after the section title
\usepackage{iftex} % check if engine is pdflatex, xetex or luatex

% Ensure that generate pdf is machine readable/ATS parsable:
\ifPDFTeX
    \input{glyphtounicode}
    \pdfgentounicode=1
    \usepackage[T1]{fontenc}
    \usepackage[utf8]{inputenc}
    \usepackage{lmodern}
\fi

\usepackage{charter}

% Some settings:
\raggedright
\AtBeginEnvironment{adjustwidth}{\partopsep0pt} % remove space before adjustwidth environment
\pagestyle{empty} % no header or footer
\setcounter{secnumdepth}{0} % no section numbering
\setlength{\parindent}{0pt} % no indentation
\setlength{\topskip}{0pt} % no top skip
\setlength{\columnsep}{0.15cm} % set column seperation
\pagenumbering{gobble} % no page numbering

\titleformat{\section}{\needspace{4\baselineskip}\bfseries\large}{}{0pt}{}[\vspace{1pt}\titlerule]

\titlespacing{\section}{
    % left space:
    -1pt
}{
    % top space:
    0.3 cm
}{
    % bottom space:
    0.2 cm
} % section title spacing

\renewcommand\labelitemi{$\vcenter{\hbox{\small$\bullet$}}$} % custom bullet points
\newenvironment{highlights}{
    \begin{itemize}[
        topsep=0.10 cm,
        parsep=0.10 cm,
        partopsep=0pt,
        itemsep=0pt,
        leftmargin=0 cm + 10pt
    ]
}{
    \end{itemize}
} % new environment for highlights


\newenvironment{highlightsforbulletentries}{
    \begin{itemize}[
        topsep=0.10 cm,
        parsep=0.10 cm,
        partopsep=0pt,
        itemsep=0pt,
        leftmargin=10pt
    ]
}{
    \end{itemize}
} % new environment for highlights for bullet entries

\newenvironment{onecolentry}{
    \begin{adjustwidth}{
        0 cm + 0.00001 cm
    }{
        0 cm + 0.00001 cm
    }
}{
    \end{adjustwidth}
} % new environment for one column entries

\newenvironment{twocolentry}[2][]{
    \onecolentry
    \def\secondColumn{#2}
    \setcolumnwidth{\fill, 4.5 cm}
    \begin{paracol}{2}
}{
    \switchcolumn \raggedleft \secondColumn
    \end{paracol}
    \endonecolentry
} % new environment for two column entries

\newenvironment{threecolentry}[3][]{
    \onecolentry
    \def\thirdColumn{#3}
    \setcolumnwidth{, \fill, 4.5 cm}
    \begin{paracol}{3}
    {\raggedright #2} \switchcolumn
}{
    \switchcolumn \raggedleft \thirdColumn
    \end{paracol}
    \endonecolentry
} % new environment for three column entries

\newenvironment{header}{
    \setlength{\topsep}{0pt}\par\kern\topsep\centering\linespread{1.5}
}{
    \par\kern\topsep
} % new environment for the header

\newcommand{\placelastupdatedtext}{% \placetextbox{<horizontal pos>}{<vertical pos>}{<stuff>}
  \AddToShipoutPictureFG*{% Add <stuff> to current page foreground
    \put(
        \LenToUnit{\paperwidth-2 cm-0 cm+0.05cm},
        \LenToUnit{\paperheight-1.0 cm}
    ){\vtop{{\null}\makebox[0pt][c]{
        \small\color{gray}\textit{Last updated in September 2024}\hspace{\widthof{Last updated in September 2024}}
    }}}%
  }%
}%

% save the original href command in a new command:
\let\hrefWithoutArrow\href

% new command for external links:


\begin{document}
    \newcommand{\AND}{\unskip
        \cleaders\copy\ANDbox\hskip\wd\ANDbox
        \ignorespaces
    }
    \newsavebox\ANDbox
    \sbox\ANDbox{$|$}

    \begin{header}
        \fontsize{25 pt}{25 pt}\selectfont Karteek Gandiboyina

        \vspace{5 pt}

        \normalsize
        \mbox{Champaign, IL}%
        \kern 5.0 pt%
        \AND%
        \kern 5.0 pt%
        \mbox{\hrefWithoutArrow{mailto:gkarteek99@gmail.com}{gkarteek99@gmail.com}}%
        \kern 5.0 pt%
        \AND%
        \kern 5.0 pt%
        \mbox{\hrefWithoutArrow{tel:+1-447-902-6750}{+1-4479026750}}%
        \kern 5.0 pt%
        \AND%
        \kern 5.0 pt%
        \mbox{Portfolio: \hrefWithoutArrow{https://muralikarteek7.github.io/}{muralikarteek7.github.io}}%
        \kern 5.0 pt%
        \AND%
        \kern 5.0 pt%
        \mbox{Github: \hrefWithoutArrow{https://github.com/muralikarteek7}{muralikarteek7}}%
    \end{header}

    \vspace{\dimexpr 5pt - 0.3cm \relax}


    \section{Summary}

        \begin{onecolentry}
            Robotics \& Computer Vision Engineer with 4+ years of industry experience building production-grade perception, SLAM, and manipulation systems. Deep expertise in 3D Gaussian Splatting, 6-DoF pose estimation, real-time object detection, and ROS-based autonomy stacks. Proven track record: 1st place in the ENACT Challenge (CVPR 2026), 2 patents filed (USPTO \& JPO), IROS 2026 publication, and industry collaborations with JAXA and Georgia Tech. MS in Autonomy \& Robotics (GPA 3.89), UIUC.
        \end{onecolentry}


    \section{Skills}

        \begin{onecolentry}
            \textbf{Computer Vision \& Perception:} 3D Gaussian Splatting (3DGS), NeRF, SLAM, 6-DoF Pose Estimation, Object Detection (Detectron2, YOLOv7/v8), Point Cloud Processing, Optical Flow, Structure-from-Motion, SAM/SAM2, Visual Odometry
            \newline
            \textbf{Robotics \& Autonomy:} ROS/ROS2, MPC, PID, Min-Snap Trajectory Planning, VIO (Basalt), Quadrotor Control, Robotic Manipulation, Bin Picking, AirSim, MuJoCo, MetaWorld, Gym AI
            \newline
            \textbf{Machine Learning:} PyTorch, JAX, CNNs, Transformers, VAEs, Diffusion Models, GNNs, World Models (DreamerV3), Actor-Critic RL, SAC, Imitation Learning (DAgger), LLMs, VLMs, VLAs, Foundational Models
            \newline
            \textbf{Languages \& Tools:} Python, C++, Docker, Linux, Git, VOXL 2
        \end{onecolentry}



    
    \section{Experience}

    \begin{twocolentry}{
            Feb 2026 – Present
        }
            \textbf{Robotics Engineer} | Rational CyPhy Inc\end{twocolentry}

        \vspace{0.05 cm}
        \begin{onecolentry}
            \begin{highlights}
                \item Developing a UAV autonomy stack on VOXL 2 + Pixhawk for wildfire-suppression BVLOS operations, integrating PX4 flight control with onboard perception for reliable flight in GPS-denied, high-interference environments.
                \item Fused GPS, IMU, and Basalt VIO into PX4 EKF2 for robust state estimation; tuned failover logic to maintain continuous EKF operation under GPS dropout, ensuring stable flight in degraded conditions.
                \item Developing a synchronized aerial data collection pipeline on VOXL 2 + Pixhawk, capturing time-aligned camera and IMU streams for offline 3DGS reconstruction targeting high-fidelity scene mapping for wildfire-suppression mission planning.
            \end{highlights}
        \end{onecolentry}


        \begin{twocolentry}{
            Sept 2025 – Jan 2026
        }
            \textbf{Graduate Research Assistant} | Prof. Sayan Mitra | Coordinated Science Laboratory at UIUC\end{twocolentry}

        \vspace{0.05 cm}
        \begin{onecolentry}
            \begin{highlights}
                \item Trained a DreamerV3 World Model on Crazyflie FPV data to learn a compact latent space where, given the current frame and action, the model predicts the next latent state and future observations; used the learned world model for imagination-based policy training without direct environment interaction.
                \item Built a 6-DoF pose estimation pipeline using 3DGS-generated synthetic FPV datasets (10k frames) and a fine-tuned U-Net, achieving sub-centimeter positional accuracy in sim-to-real transfer.
                \item Developed an actor-critic policy (using PPO) within the latent world model for vision-based autonomous navigation, increasing successful waypoint completion from 70 \% to 95 \% in simulation trials
                \item Designed min-snap trajectory planners with velocity and acceleration limits using differential flatness, reducing planning time by 30 \% and enabling agile quadrotor maneuvers in cluttered environments
            \end{highlights}
        \end{onecolentry}


        \vspace{0.1 cm}


        
        \begin{twocolentry}{
            July 2021 – Aug 2024
        }
            \textbf{R\&D Robotics Engineer} | Konica Minolta -- Tokyo, Japan\end{twocolentry}

        \vspace{0.05 cm}
        \begin{onecolentry}
            \begin{highlights}
                \item Engineered a structured-light point cloud sensing device achieving \textbf{0.99 IoU} on 1\,cm\textsuperscript{2} items at 0.05\,m range under $<$10\,Lux enabling bin picking of highly reflective industrial components.
                \item Built an auto-annotation pipeline (SAM,Detectron2, YOLOv7) that reduced manual labeling time by $>$80\%, accelerating object detection model training cycles across multiple product lines.
                \item Developed a language-conditioned multi-task RL framework (SAC + DAgger) with Georgia Tech; improved zero-shot task success rates by \textbf{200\%} and enabled rapid skill transfer via natural language goal specification.
                \item Collaborated with JAXA to design a vision-based 6-DoF grasping system for multi-limbed robots aboard the International Space Station, optimized for microgravity manipulation.
                \item Filed \textbf{2 patents} in computer vision and robotic grasping (USPTO \& JPO).
            \end{highlights}
        \end{onecolentry}


        \vspace{0.1 cm}

        \begin{twocolentry}{
            April 2020 – July 2020
        }
            \textbf{Machine Learning Intern} | Philips Innovation Campus -- Bangalore, India\end{twocolentry}

        \vspace{0.05 cm}
        \begin{onecolentry}
            \begin{highlights}
                \item Developed a 3D residual U-Net region-proposal network for pulmonary nodule detection on 888 CT scans (LUNA16); achieved a \textbf{FROC score of 0.914}, surpassing the baseline by 8\%.
            \end{highlights}
        \end{onecolentry}




    
    \section{Projects}

        \begin{twocolentry}{
            \href{https://github.com/muralikarteek7/DroneRacer-MPC-Vision.git}{DroneRacer-MPC-Vision.git}
        }
           {\textbf{Autonomous Drone Racing} | UIUC}\end{twocolentry}

        \vspace{0.05 cm}
        \begin{onecolentry}
            \begin{highlights}
                \item Implemented MPC \& PID spline tracker for autonomous gate navigation; achieved a Tier-1 race record of \textbf{50.13 seconds}.
                \item Deployed NanoSAM with a keypoint detector for misaligned gate correction, reaching a positional prediction error of \textbf{0.05\,m}. \textit{Stack: Python, NeRF, AirSim, MPC, PID, Trajectory Planning, State-Space Estimation.}
            \end{highlights}
        \end{onecolentry}

        \vspace{0.1 cm}

        \begin{twocolentry}{
            \href{https://github.com/muralikarteek7/VLM4Autonomy-.git}{VLM4Autonomy-.git}
        }
            {\textbf{VLM4Autonomy} | UIUC}\end{twocolentry}

        \vspace{0.05 cm}
        \begin{onecolentry}
            \begin{highlights}
                \item Fused SAM2, optical flow, and VLMs for real-time object tracking and ego-vehicle motion estimation; integrated YOLOv8 and Structure-from-Motion for spatial grounding. \textit{Stack: Python, VLM, SAM2, YOLOv8, Visual Odometry, SpatialBot.}
            \end{highlights}
        \end{onecolentry}

        \vspace{0.1 cm}

        \begin{twocolentry}{
            \href{https://github.com/muralikarteek7/Transformer-VAE-Stock-Predictor/tree/main}{VAE-Stock-Predictor.git}
        }
            {\textbf{Volatility-Aware Stock Prediction via Transformer-VAE-Flow} | UIUC}\end{twocolentry}

        \vspace{0.05 cm}
        \begin{onecolentry}
            \begin{highlights}
                \item Built a Transformer-VAE with RealNVP normalizing flows conditioned on VIX for synthetic financial time-series generation; achieved \textbf{Wasserstein Distance of 0.0629}. \textit{Stack: PyTorch, VAE, RealNVP, Transformers.}
            \end{highlights}
        \end{onecolentry}
        


        


   \section{Awards \& Honors}

        \begin{twocolentry}{
            May 2026
        }
            \textbf{1st Place, ENACT Embodied Cognition Challenge} | CVPR 2026 FMEA Workshop\end{twocolentry}

        \vspace{0.05 cm}
        \begin{onecolentry}
            \begin{highlights}
                \item Fine-tuned Qwen2.5-VL-7B with LoRA SFT on 7,672 samples; achieved \textbf{94.6\% Task Accuracy} and \textbf{98.4\% Pairwise Accuracy}, outperforming GPT-5 and Gemini 2.5 Pro zero-shot on the hidden test set. \href{https://github.com/muralikarteek7/ENACT-Qwen2.5VL}{[GitHub]}
                \item Invited to give a \textbf{5-minute Contributed Talk} at the CVPR 2026 FMEA Workshop (June 2026).
            \end{highlights}
        \end{onecolentry}

        \vspace{0.1 cm}

   \section{Patents \& Publications}

        \begin{twocolentry}{
            \href{https://patents.google.com/patent/US20250336175A1/en}{IROS 2026}
        }
            \textbf{FalconTrack: Photorealistic Auto-Labeled Perception and Physics-Aware Vision-Based Aerial Tracking}\end{twocolentry}

        \vspace{0.05 cm}
        \begin{onecolentry}
            \begin{highlights}
                \item Developed a unified 3DGS-based simulator that automates 6-DoF pose and mask labeling (10k images in $<$20 min), enabling zero-shot sim-to-real aerial tracking with \textbf{100\% success} on high-speed trajectories.
            \end{highlights}
        \end{onecolentry}

        \vspace{0.1 cm}

        \begin{twocolentry}{
            \href{https://patents.google.com/patent/US20250336175A1/en}{US: 20250336175}
        }
            \textbf{Learning Data Generation Support Device, Movement Controller, and Article Acquisition System}\end{twocolentry}

        \vspace{0.05 cm}
        \begin{onecolentry}
            \begin{highlights}
                \item Co-invented an automated data generation framework to eliminate manual labeling bottlenecks for complex 3D pick-and-place robotics in bulk manufacturing.
            \end{highlights}
        \end{onecolentry}

        \vspace{0.1 cm}

        \begin{twocolentry}{
            \href{https://patents.justia.com/patent/20240338849}{JPO: 2023-060687}
        }
            \textbf{Component Posture Information Acquiring Device and Posture Determination Method}\end{twocolentry}

        \vspace{0.05 cm}
        \begin{onecolentry}
            \begin{highlights}
                \item Invented a vision-based sensing system using ROI determination and feature extraction for accurate component posture estimation during robotic pick-and-place.
            \end{highlights}
        \end{onecolentry}


    \section{Education}

        \begin{twocolentry}{Aug 2024 -- Dec 2025}
            {\textbf{University of Illinois Urbana-Champaign} | MS in Autonomy \& Robotics | GPA: 3.89/4.0}\end{twocolentry}

        \vspace{0.05 cm}
        \begin{onecolentry}
            \begin{highlights}
                \item \textbf{Coursework:} Deep Learning for Graphs, Computer Vision, Safe Autonomy, Deep Generative Models
            \end{highlights}
        \end{onecolentry}

        \vspace{0.05 cm}

        \begin{twocolentry}{Aug 2017 -- May 2021}
            \textbf{Indian Institute of Technology Kharagpur} | B.Tech in Electrical Engineering | GPA: 3.48/4.0\end{twocolentry}

        \vspace{0.05 cm}
        \begin{onecolentry}
            \begin{highlights}
                \item \textbf{Coursework:} Deep Learning, Machine Learning, Embedded Systems, Control Systems
            \end{highlights}
        \end{onecolentry}


\end{document}